%%%%%%%%%%  scan idx 355  |  folio 342  %%%%%%%%%%
\SGAPrintPage{342}%
sending \((1,0)\) to \(1\), and \((0,x)\) to \(x\mathbin{.}\lambda_{-1}(N)\). It follows immediately from (5.5.1) and (5.3.1) that this is a \(\lambda\)-homomorphism, and, since \(A\) is an integral domain in the case considered, one immediately verifies that it is injective; since \(A\) is a \(\lambda\)-ring, we deduce the result.

In the \(\lambda\)-ring \(A_N\), we can define operations \(\gamma\) (3.2). For \(x\in A\), the value of \(\gamma^p\) calculated on the element \((0,x)\) of \(A_N\) is given by:
\[
  \gamma^p(0,x)=\sum_{i=0}^{p-1}\binom{p-1}{i}\lambda^{p-i}(N,x),
\]
it is therefore an element of \(A\); it is defined by the relation:
\sgaStableTarget{sga6.e05.eq.5-5-2}\begin{equation}\tag{5.5.4}
  \gamma^p(N,x)\mathbin{.}\lambda_{-1}(N)=\gamma^p(x\mathbin{.}\lambda_{-1}(N)),
\end{equation}
and the universal polynomials resulting from it.

\sgaStableTarget{sga6.e05.para.5-6}5.6. Let \(A\) be the \(\lambda\)-ring generated over \(\mathbb Z\) by the generators \(N,L,x\), subject to the relations \(\lambda^k(N)=0\) for \(k>d\), \(\lambda^k(L)=0\) for \(k>1\). We propose to calculate the \(\lambda^p(N+L,x)\) in terms of the \(\lambda^q(N,x)\).

By definition:
\[
  \lambda^p(N+L,x)\mathbin{.}\lambda_{-1}(N)\mathbin{.}\lambda_{-1}(L)
  =\lambda^p(x\mathbin{.}\lambda_{-1}(N)\mathbin{.}\lambda_{-1}(L)),
\]
that is:
\[
\begin{aligned}
1&+\lambda_{-1}(N)\mathbin{.}\lambda_{-1}(L)\mathbin{.}\sum_{k\ge1}\lambda^k(N+L,x)\mathbin{.} t^k \\
 &=\lambda_t(x\mathbin{.}\lambda_{-1}(N))\circ \lambda_t(1-L) \\
 &=\frac{\lambda_t(x\mathbin{.}\lambda_{-1}(N))}{\lambda_{Lt}(x\mathbin{.}\lambda_{-1}(N))}
   \qquad\left(\text{since }\lambda_t(1-L)=\frac{\lambda_t(1)}{\lambda_t(L)}=\frac{1+t}{1+Lt}\text{, and we apply }(2.3.2)\right) \\
 &=\frac{1+\ldots+\lambda^k(x\mathbin{.}\lambda_{-1}(N))t^k+\ldots}{1+\ldots+\lambda^k(x\mathbin{.}\lambda_{-1}(N))L^k t^k+\ldots} \\
 &=1+\frac{\sum_{k\ge1}\lambda^k(N,x)\mathbin{.}\lambda_{-1}(N)\mathbin{.}(1-L^k)t^k}{1+\sum_{k\ge1}\lambda^k(N,x)\mathbin{.}\lambda_{-1}(N)\mathbin{.} L^k t^k}
\end{aligned}
\]

%%%%%%%%%%  scan idx 356  |  folio 343  %%%%%%%%%%
\SGAPrintPage{343}%
Thus \(\lambda^p(N+L,x)\) is the coefficient of \(t^p\) in the series:
\sgaStableTarget{sga6.e05.eq.5-6-1}\[
  \frac{\displaystyle\sum_{k\ge1}\lambda^k(N,x)
        (1+\ldots+L^{k-1})t^k}
       {\displaystyle 1+\sum_{k\ge1}\lambda^k(N,x)
        \lambda_{-1}(N)L^k t^k}.
\]

\sgaStableTarget{sga6.e05.prop.5-7}\textbf{Proposition 5.7.} \uline{Let} \(K\) \uline{be a binomial ring}, \(A\) \uline{an augmented} \(K\)-\(\lambda\)-\uline{algebra}, \uline{and} \(N\) \uline{an element of} \(A\) \uline{such that} \(\lambda^k(N)=0\) \uline{for} \(k>d\). \uline{Then for every} \(x\in A\) \uline{and every} \(p\ge1\), \uline{we have}:
\sgaStableTarget{sga6.e05.eq.5-7-1}\begin{equation}\tag{5.7.1}
  \gamma^p(N,x)\in\Fil^{p-d}A.
\end{equation}

Let \(A_0\) be the \(K\)-\(\lambda\)-algebra generated by generators \(N'_0,X'_0\), subject to the relations \(\gamma^k(N'_0)=0\) for \(k>d\). We define the augmentation homomorphism \(A_0\xrightarrow{\epsilon}K\) by \(\epsilon(N'_0)=\epsilon(X'_0)=0\). We know that \(A_0\) is the ring of polynomials with coefficients in \(K\) in the indeterminates \(\gamma^i(N'_0)\), with \(1\le i\le d\), and \(\gamma^j(X'_0)\), with \(j\ge1\), and that the \(\lambda\)-filtration on \(A_0\) is identical to the weight filtration, assigning the weight \(i\) to \(\gamma^i(N'_0)\), \(\gamma^i(X'_0)\) (4.12 and 4.15).

Put \(N_0=N'_0+d\), \(X_0=X'_0+\epsilon(x)\). We shall first prove 5.7 for \(N_0\) and \(X_0\). By the preceding discussion, (5.7.1) is equivalent to:
\[
  \gamma^p(N_0,X_0)\mathbin{.}\gamma^d(N'_0)\in\Fil^pA_0.
\]
Now \(\gamma^d(N'_0)=(-1)^d\lambda_{-1}(N_0)\); this relation can therefore be written:
\[
  (-1)^d\gamma^p(X_0\gamma^d(N_0))\in\Fil^pA_0
\]
by (5.5.4). This relation does hold, since \(\gamma^d(N_0)\) has zero augmentation.

By 4.7, there exists a \(K\)-\(\lambda\)-homomorphism \(A_0\longrightarrow A\) compatible with the augmentations, sending \(N'_0\) to \(N-d\), and \(X'_0\) to \(x-\epsilon(x)\), hence \(N_0\) to \(N\) and \(X_0\)
%%%%%%%%%%  scan idx 357  |  folio 344  %%%%%%%%%%
\SGAPrintPage{344}%
to \(x\), whence the result, since the homomorphism is compatible with the filtrations.

\sgaStableTarget{sga6.e05.sec.6}\section*{6. Chern ring}

In this section, the base ring will be a binomial ring (2.7) \(K\), fixed once and for all. We shall use the category \(\mathcal C\) (1.1), and therefore return to the notation of Section 1, reserving the notation \(\widehat G\) for the functor defined on \(\mathcal C\) (1.1), and denoting by \(\widehat G\circ i\) its restriction to \(\mathcal C_0\) (1.8).

\sgaStableTarget{sga6.e05.para.6-1}6.1. In 3.5 and 3.6 we defined morphisms of functors:
\[
  (\widehat G\circ i)\times(\widehat G\circ i)\xrightarrow{\nabla}\widehat G\circ i,
  \qquad
  \widehat G^t\circ i\xrightarrow{\gamma'}(\widehat G^u\circ i)\circ(\widehat G^t\circ i)_\nabla,
  \qquad
  \widehat G\circ i\xrightarrow{a}\widehat G\circ i,
\]
the last being multiplication under the law \(\nabla\) by \(\gamma_t(a)\), with \(a\in K\).

By 1.9, we deduce morphisms:
\[
  \widehat G\times\widehat G\xrightarrow{\nabla}\widehat G_g,
  \qquad
  \widehat G^t\xrightarrow{\gamma'}\widehat G^u\circ i\circ\widehat G^t_{g\nabla},
  \qquad
  \widehat G\xrightarrow{a}\widehat G_g;
\]
and it follows from the isobaricity properties of the polynomials \(P'_n\), \(Q'_{i,j}\), \(R_{n,a}\) (3.5, 3.6, 2.9) that the image of the first and third functors is in \(\widehat G\), and that of the second in the subgroup of formal series in \(u\) whose coefficients in nonzero degree are in \(\widehat G^t\).

Composing the first and third with the canonical morphism: \(\widehat G\longrightarrow 1+\widehat{\operatorname{Gr}\,\mathord{.}\,}^{+}\) (1.11), we obtain, by 1.12, two morphisms of functors on \(\mathcal G\):
%%%%%%%%%%  scan idx 358  |  folio 345  %%%%%%%%%%
\SGAPrintPage{345}%
\[
  (1+\widehat{\mathord{.}}{}^{+})\times(1+\widehat{\mathord{.}}{}^{+})\xrightarrow{\ \ast\ }1+\widehat{\mathord{.}}{}^{+},
\]
\[
  1+\widehat{\mathord{.}}{}^{+}\xrightarrow{\ a\ }1+\widehat{\mathord{.}}{}^{+},
\]
defining on the functor \(1+\widehat{\mathord{.}}{}^{+}\) a multiplication and operations of \(K\). By 1.13, the multiplication is defined by its value on the element \((1+X,1+Y)\) of \((1+\widehat{K[X,Y]}{}^{+})^{2}\); it will be denoted by \(\ast\); we therefore have:
\begin{equation}\tag{6.1.1}
  (1+X)\ast(1+Y)=\frac{1+X+Y}{(1+X)(1+Y)}.
\end{equation}
The \(n\)-th universal polynomial of this law is the isobaric component of weight \(n\) of \(P'_n\) (1.15), and will be denoted by \(P''_n\).

The operations of \(K\) will again be denoted by exponentiation (3.8). The universal polynomials defining these operations are the isobaric components of lowest degree of the \(R_{n,a}\) (1.15); since these are isobaric, they are the \(R_{n,a}\) themselves.

Let \(A\in\operatorname{Ob}(\mathcal G)\); the ring \(1+\widehat A^{+}\) defined by the preceding laws is in general without a unit; we therefore extend this ring by adjoining the scalars of \(K\), introducing the new ring \(K\times(1+\widehat A^{+})\), endowed with the composition laws defined by: \(\forall\,a,a'\in K\), \(\varphi,\varphi'\in1+\widehat A^{+}\),
\[
  (a,\varphi)+(a',\varphi')=(a+a',\varphi\varphi')~;
\]
\begin{equation}\tag{6.1.2}
  (a,\varphi).(a',\varphi')=(aa',\varphi^{a'}\mathbin{.}{\varphi'}^{\,a}\mathbin{.}(\varphi\ast\varphi')).
\end{equation}

The ring thus constructed admits (1,1) as its unit, and \(1+\widehat A^{+}\) is identified with the ideal \((0)\times(1+\widehat A^{+})\). It is naturally endowed with a \(K\)-algebra structure such that:
\begin{equation}\tag{6.1.3}
  a.(a',\varphi)=(a,1).(a',\varphi)=(aa',\varphi^{a}).
\end{equation}

%%%%%%%%%%  scan idx 359  |  folio 346  %%%%%%%%%%
\SGAPrintPage{346}%
We endow it with a \(K\)-\(\lambda\)-algebra structure as follows: we observed above that the image of the functor \(\gamma'\colon\widehat G^t\longrightarrow\widehat G^u\circ i\circ\widehat G^t_g\) is in the subgroup of formal series in \(u\) whose coefficients in nonzero degree are in \(\widehat G^t\). Using the canonical morphism: \(\widehat G\longrightarrow 1+\widehat{\operatorname{Gr}\,\mathord{.}\,}^{+}\), which gives: \(\widehat G\longrightarrow K\times(1+\widehat{\operatorname{Gr}\,\mathord{.}\,}^{+})\), we finally deduce a morphism:
\[
  \widehat G^t\xrightarrow{\gamma'}\widehat G^u\circ i\circ(K\times(1+\widehat{\operatorname{Gr}\,\mathord{.}\,}^{+})),
\]
which, by 1.12, defines a morphism of functors on \(\mathcal G\):
\[
  1+\widehat{\mathord{.}}{}^{+}\xrightarrow{\gamma'}\widehat G^u\circ i\circ(K\times(1+\widehat{\mathord{.}}{}^{+})).
\]
The universal polynomials defining this morphism are the coefficients of the \({\gamma'}^{i}(1+X_1+\ldots+X_n+\ldots)\); the \(j\)-th coefficient is the isobaric component of weight \(j\) of \(Q'_{i,j}\) (a proof analogous to that of 1.15), and will be denoted by \(Q''_{i,j}\).

Finally, we extend the morphism thus constructed to \(K\times(1+\widehat{\mathord{.}}{}^{+})\) naturally, using the \(\lambda\)-structure of \(K\). We therefore obtain a pre-\(\lambda\)-ring structure on \(K\times(1+\widehat{\mathord{.}}{}^{+})\), defined by the corresponding operations \(\gamma\), which satisfy:
\begin{equation}\tag{6.1.4}
  \gamma_t(a,\varphi)=\gamma_t(a).\Bigl(1+\sum_{p\ge1}{\gamma'}^{p}(\varphi)t^{p}\Bigr),
\end{equation}
and which will be denoted by \(\gamma\) instead of \(\gamma'\), since no confusion is possible.

To verify that this is a \(\lambda\)-ring structure, we verify the commutativity of (3.7.2); it is known for the elements of \(K\); for the elements of \(1+\widehat{\mathord{.}}{}^{+}\), it follows from that of (3.7.2) for \(\widehat G\) by passage to the quotient. The result follows at once.

\providecommand{\Ghat}{\widehat G}
\providecommand{\Aqhat}{\widehat A_{\mathbb Q}}
\providecommand{\Ch}{\operatorname{Ch}}
\providecommand{\Log}{\operatorname{Log}}
\providecommand{\id}{\operatorname{id}}
\providecommand{\gr}{\operatorname{gr}}
\providecommand{\Hom}{\operatorname{Hom}}
\providecommand{\epsi}{\varepsilon}
\providecommand{\clgr}{\operatorname{cl}}
\providecommand{\Fil}{\operatorname{Fil}}
\providecommand{\Gr}{\operatorname{Gr}}
\providecommand{\Sym}{\operatorname{Sym}}

% Local macros for Exposé V, §§6--7
\providecommand{\Ghat}{\widehat G}
\providecommand{\Aqhat}{\widehat A_{\mathbb Q}}
\providecommand{\Ch}{\operatorname{Ch}}
\providecommand{\Log}{\operatorname{Log}}
\renewcommand{\id}{\operatorname{id}}
\providecommand{\gr}{\operatorname{gr}}
\providecommand{\Hom}{\operatorname{Hom}}
\providecommand{\epsi}{\varepsilon}
\providecommand{\clgr}{\operatorname{cl}}
\providecommand{\Fil}{\operatorname{Fil}}
\providecommand{\Gr}{\operatorname{Gr}}
\providecommand{\Sym}{\operatorname{Sym}}

